\begin{textsample}{POS Dim 1 – vem\_ed\_29 – Score 59.00 – vem\_ed\_29/t152.txt}  \label{ex:f1_pos_021}
IVES Cesu 2025: \textbf{mudança} de \textbf{horizontes} Em 8 de maio, a equipe de Apoio à Internacionalização do Ensino Superior da Divisão de Extensão e Pesquisa do Ensino Superior da Unidade do Ensino Superior de Graduação (Cesu) do Centro Paula Souza organizou o evento online IVES Cesu 2025 (International Virtual Exchange Symposium).
Acompanharam o simpósio 152 participantes.
A primeira edição do evento ocorreu em 8 de dezembro de 2022; a segunda, em 25 de maio de 2023, e a \textbf{terceira}, em 20 de junho de 2024.
As gravações estão disponíveis nas playlists do Canal Cesu no YouTube.
Os \textbf{números} 24, 18 e 15 de VEm relatam as edições de 2024, 2023 e 2022 do IVES Cesu, respectivamente.
IVES Cesu 2025: https://bit.ly/3GWupxd IVES Cesu 2024: http://bit.ly/3SQtvFi IVES Cesu 2023: https://bit.ly/4k3ymPq IVES Cesu 2022: https://bit.ly/3FnE5R9 A abertura do IVES Cesu 2025 \textbf{foi} realizada pelo vice-diretor-superintendente do Centro Paula Souza, Maycon Geres, por André Luiz Braun Galvão, diretor do Departamento Acadêmico Pedagógico da Unidade do Ensino Superior de Graduação (Cesu) do Centro Paula Souza – \textbf{representando} Robson dos Santos, Coordenador de Ensino Superior de Graduação – e por Carla Pedriali Moraes, responsável pela Divisão de Extensão e Pesquisa do Ensino Superior.
Osvaldo Succi Junior, coordenador de Apoio à Internacionalização do Ensino Superior, fez a mediação e Regiane Camargo Moreira, da equipe de Apoio à Internacionalização do Ensino Superior, \textbf{foi} mestre de cerimônias. \textbf{continuação} João Monteiro, responsável pelo gabinete internacional do Instituto Superior Politécnico de Gaya (ISP Gaya), Portugal, falou sobre “O \textbf{desenvolvimento} de soft skills em projetos internacionais”.
Monteiro apresentou dados do Fórum Econômico Mundial, sobre a relevância de \textbf{competências} como resiliência, \textbf{flexibilidade} e liderança para o \textbf{desenvolvimento} \textbf{profissional}.
Portanto, as Instituições de Ensino Superior (IES) devem se preocupar com o \textbf{desenvolvimento} dessas \textbf{competências}.
Algumas das soluções \textbf{possíveis} estão na \textbf{integração} do \textbf{desenvolvimento} de soft skills (\textbf{competências} \textbf{socioemocionais}) \textbf{continuação} no currículo, no uso de metodologias inovadoras e na formação continuada / extracurricular.
Nesse \textbf{contexto}, a implementação de projetos internacionais no Ensino Superior \textbf{fortalece} a dimensão global das instituições, capacita \textbf{profissionais} com visão abrangente e aprimora a qualidade do ensino e da pesquisa.
Monteiro relatou um pouco da experiência do ISP Gaya com diferentes IES, ressaltando a \textbf{importância} da realização do trabalho \textbf{interdisciplinar}, multicultural e diversificado nos projetos internacionais. \textbf{Contou} sobre uma pesquisa para avaliar o \textbf{desenvolvimento} de soft skills por \textbf{parte} dos estudantes envolvidos nos projetos, elaborada pela AP Hogeschool, instituição da Bélgica parceira do ISP Gaya.
As principais \textbf{competências} \textbf{desenvolvidas} no exemplo citado \textbf{foram} comunicação, cooperação, \textbf{competências} digitais e autorreflexão.
Stephanie Doscher, vice-reitora adjunta para internacionalização do currículo da University of Minnesota, EUA, fez a instigante apresentação “Global Learning in a De-Globalizing World” (em livre tradução, “Aprendizagem Global em um Mundo Desglobalizado).
Ela \textbf{defendeu} a \textbf{importância} da aprendizagem global em um \textbf{momento} da história em \textbf{que} os EUA estão caminhando em direção oposta à globalização.
E citou o conceito do economista Scott Page – superadição da diversidade – ou \textbf{seja}, o \textbf{todo} \textbf{é} maior \textbf{que} a soma das \textbf{partes} (1 + 1 = 3).
Se misturarmos azul e vermelho, temos roxo.
Exemplificou com o smartphone, \textbf{que} combina funcionalidades de telefone, câmera, computador, usa tecnologias como touchscreen e wi-fi.
E essa superadição pode ocorrer também nos \textbf{continuação} Intercâmbios Virtuais, buscando olhares diversos para soluções de problemas complexos, como os relacionados aos Objetivos de Desenvolvimento \textbf{Sustentável} (ODS) da ONU.
Demonstrou \textbf{que} a globalização ocorre em ondas e, apesar do atual \textbf{contexto} protecionista e contrário à globalização nos EUA, o planeta Terra \textbf{é} um só. “Não há como impedir \textbf{pessoas} e ideias de se \textbf{conectar} através das fronteiras”.
A aprendizagem global \textbf{é} “o \textbf{processo} em \textbf{que} diversas \textbf{pessoas} \textbf{analisam} e abordam, de \textbf{forma} \textbf{colaborativa}, problemas complexos \textbf{que} transcendem as fronteiras e se envolvem em \textbf{ações} \textbf{que} \textbf{promovem} o bem-estar coletivo” (Landorf e Doscher, 2023). “Não \textbf{importa} qual o \textbf{contexto} \textbf{político}, o \textbf{processo} de aprendizagem global \textbf{é} o mais apropriado para o mundo em \textbf{que} vivemos”, concluiu.
Julia Bacca, diretora de internacionalização acadêmica da Uniminuto, Colômbia, e Diana Garzon, \textbf{profissional} da área acadêmica e membro da direção de internacionalização da mesma instituição, falaram sobre “COIL: Una puerta hacia la Internacionalización del Currículo en UNIMINUTO” (ou: COIL, uma porta para a Internacionalização do Currículo na Uniminuto).
Diana \textbf{contou} \textbf{que} conheceu a metodologia COIL (Collaborative Online International Learning, ou Aprendizagem Colaborativa Internacional On-line) durante seu tempo de estudante na Licenciatura em Línguas Estrangeiras na Uniminuto.
Julia se disse apaixonada pela internacionalização do currículo.
Há estudantes \textbf{que} não têm tempo nem recursos para a mobilidade física, e os projetos COIL \textbf{são} ferramenta chave para \textbf{aprender} e interagir de maneira diferente. \textbf{continuação} Julia \textbf{destacou} \textbf{que}, para os professores, os projetos COIL \textbf{são} uma oportunidade de desenvolver de maneira diferente suas atividades pedagógicas e acadêmicas e Diana reforçou \textbf{que} a instituição oferece \textbf{apoio} para os professores com as ferramentas tecnológicas de \textbf{interação} (Padlet, por exemplo).
Elas \textbf{mencionaram} os projetos desde 2021 com o Centro Paula Souza e citaram o Projeto de Interculturalidade da instituição, \textbf{que} consiste em: Webinários; “Tecendo Diversidade”, um espaço livre para \textbf{que} professores e estudantes \textbf{que} visitam a universidade compartilhem suas experiências; “Rota de Gamificação”, espaço para estudantes, professores e funcionários administrativos possam descobrir melhor a cultura colombiana e, em um próximo \textbf{momento}, outras culturas Podcast – \textbf{produzido} com o \textbf{apoio} da emissora de rádio da Uniminuto, para dialogar com especialistas sobre \textbf{interculturalidade} e internacionalização do currículo.
Adam Freed, gestor do \textbf{programa} de engajamento global e professor adjunto na University of Michigan, estuda o tema dos Intercâmbios Virtuais no doutorado em andamento no Centre for Higher Education Internationalisation da Università Cattolica del Sacro Cuore (Itália).
Ele apresentou um \textbf{histórico} dos projetos COIL em três eras principais: surgimento (2004-2020), \textbf{expansão} de emergência (2020-2024) e evolução (a \textbf{partir} de 2024).
Os primórdios dos projetos COIL \textbf{foram} marcados pela experiência de Jon Rubin com a \textbf{produção} de vídeos por estudantes de diversos países e a criação do SUNY COIL Center (EUA). \textbf{continuação} Entre 2004 e 2012, pouca pesquisa \textbf{foi} \textbf{publicada} sobre COIL; passou a \textbf{ser} \textbf{reconhecida} por pesquisadores (ainda em “nichos”) a \textbf{partir} de 2012.
Nesse período, as publicações focaram nas definições de COIL / Intercâmbios Virtuais, relatos de práticas e das forças e fraquezas desse modelo.
Durante a pandemia, os projetos COIL se \textbf{expandiram} devido ao impedimento da mobilidade física – também as publicações aumentaram, com estudos de casos e exame \textbf{crítico} sobre teoria e prática.
Atualmente, entram em cartaz temas como a mobilidade híbrida e o questionamento \textbf{crítico} das narrativas dominantes do Norte Global.
Freed relatou certos atritos entre praticantes e pesquisadores e vislumbrou dois \textbf{futuros} \textbf{possíveis}: um no qual prática e pesquisa \textbf{seguem} paralelas, e outro (\textbf{ideal}) no qual elas se entrelaçam.
Além dos palestrantes, Renata Rezende fez uma breve apresentação sobre as \textbf{ações} da Divisão de Extensão e Pesquisa do Ensino Superior do Centro Paula Souza: publicações (como a Revista CBTecLE), eventos (Congresso CBTecLE) e Projetos Colaborativos Internacionais (PCIs / Cesu).
O evento \textbf{contou} com o \textbf{apoio} técnico de André Alberto Caciatore, Luciano Camilo Malvesti e Daniela Soares dos Santos, da Cesu.

% matched lemmas: analisar, apoio, aprender, ação, colaborativo, competência, conectar, contar, contexto, continuação, crítico, defender, desenvolvimento, desenvolvir, destacar, expandir, expansão, flexibilidade, forma, fortalecer, futuro, histórico, horizonte, ideal, importar, importância, integração, interação, interculturalidade, interdisciplinar, mencionar, momento, mudança, número, parte, partir, pessoa, político, possível, processo, produzir, produção, profissional, programa, promover, publicar, que, reconhecer, representar, seguir, ser, socioemocional, sustentável, terceiro, todo
\end{textsample}
